% ***************************************************************************************************
%
%	Szablon pracy magisterskiej dla Politechniki Wrocławskiej w wersji dwustronnej.
%	Autor:	Tomasz Strzałka
%
% ***************************************************************************************************

% Styl dwustronny z domyślną wielkością czcionki 10pt oraz oddzieloną stroną tytułową (titlepage).
% Domyślnie rodziały rozpoczynają się na stronie prawej (openright).
\documentclass[a4paper,11pt,twoside,openright]{book}
% \RequirePackage{luatex85,shellesc}
\usepackage{luatex85}
\usepackage{shellesc}

\renewcommand{\baselinestretch}{1.2} 

% ***************************************************************************************************
% Ustawienia języka
% ***************************************************************************************************

% Podstawowe ustawienia języka, według którego formatowany będzie dokument
\usepackage[english]{babel}
\usepackage{csquotes}
% Pakiet babel dla polskiego języka powoduje konflikt z pakietem amssymb.
% Polecenie '\lll' definiują oba pakiety - porządana jest druga definicja.

% W przypadku wielojęzykowości ustawia główny język dokumentu
\selectlanguage{english}

% Kodowanie dokumentu
% \usepackage[utf8]{inputenc}

% Dowolny rozmiar czcionek, kodowanie znaków
\usepackage{lmodern}

% Polskie wcięcia akapitów
% \usepackage{indentfirst}

% Polskie łamanie wyrazów
% \usepackage[plmath]{english}

% Przecinek w wyrażeniach matematycznych zamiast kropki
% \usepackage{icomma}

% Polskie formatowanie typograficzne
% \frenchspacing

% Zapewnia liczne usprawnienia wyświetlania i organizacji matematycznych formuł. 
\usepackage{amsmath}

% Wprowadza rozszerzony zestaw symboli m.in. \leadsto
\usepackage{amssymb}

% Dodatkowa, ,,kręcona'' czcionka matematyczna
\usepackage{mathrsfs}

% Dodatkowe wsparcie dla środowiska mathbb, które nie wspiera domyślnie cyfr (\mathbb{})
\usepackage{bbold}


% Fixes/improves amsmath
\usepackage{mathtools}


% ***************************************************************************************************
% Kolory  
% ***************************************************************************************************

% Umożliwia kolorowanie poszczególnych komórek tabeli
\usepackage[table]{xcolor}% http://ctan.org/pkg/

% Umożliwia łatwą zmianę koloru linii w tabeli
\usepackage{tabu}

% Umożliwia rozszerzoną kontrolę nad kolorami.
\usepackage{xcolor}

% Definicje kolorów
\definecolor{lgray}{HTML}{9F9F9F}
\definecolor{dgray}{HTML}{5F5F5F}
% lgray				-	nazwa nowo zdefiniowanego koloru
% HTML				-	model kolorów
% CCCCCC			-	wartość koloru zgodna z modelem

% ***************************************************************************************************
% Algorytmy 
% ***************************************************************************************************

% Udostępnia środowisko do konstruowania pseudokodów
% \usepackage[ruled,vlined,linesnumbered,longend,algochapter]{algorithm2e}
% ruled	- poziome kreski na początku i końcu algorytmu, podpis na górze oddzielony również kreską poziomą
% vlined - pionowe kreski łączące początek polecenia z jego końcem
% linesnumbered	- numerowanie kolejnych wierszy algorytmu
% longend - długie końcówki np. ifend, forend itd.
% algochapter - numeracja z rozdziałami

% Zmiana rozmiaru komentarzy
% \newcommand\algcomment[1]{
	% \footnotesize{#1}
% }

% Ustawienie zadanego stylu dla komentarzy
% \SetCommentSty{algcomment}


% \usepackage{newtxtext}
% \usepackage{newtxmath}

% Wyśrodkowana tylda
\usepackage{textcomp}%
\newcommand{\textapprox}{\raisebox{0.5ex}{\texttildelow}}

% Listowanie kodów źródłowych
\usepackage{listings} 

% ***************************************************************************************************
% Marginesy 
% ***************************************************************************************************

% Ustawienia rozmiarów stron i ich marginesów
\usepackage[headheight=18pt, top=25mm, bottom=25mm, left=25mm, right=25mm, bindingoffset=5mm]{geometry}
% headheight		-	wysokość tytułów
% top				-	margines górny
% bottom			-	margines dolny
% left				-	margines lewy
% right				-	margines prawy

% Usunięcie górnego marginesu dla środowisk
\makeatletter
\setlength\@fptop{0\p@}	
\makeatother

% ***************************************************************************************************
% Styl 
% ***************************************************************************************************

% Definiuje środowisko 'titlingpage', które zapewnia pełną kontrolę nad układem strony tytułowej.
\usepackage{titling}


% Umożliwia modyfikowanie stylu spisu treści
\usepackage{tocloft}	

\tocloftpagestyle{tableOfContentStyle}

% Definiowanie własnych stylów nagłówków i/lub stopek
\usepackage{fancyhdr}

% Domyślny styl dla pracy 
\fancypagestyle{custom}{
	\fancyhf{}									% wyczyść stopki i nagłówki
	\fancyhead[RO]{								% Prawy, nieparzysty nagłówek
		\hrulefill \hspace{16pt} \large Chapter \thechapter
		\put(-455.4, 13.5){%
			\makebox(0,0)[l]{%
				\includegraphics[width=0.05\textwidth]{Figures/pwr-logo}
			}
		}
		\put(-426,5.5){%
			\makebox(0,0)[l]{%
				\small Wrocław University of Science and Technology
			}
		}
	}
	\fancyhead[LE]{								% Lewy, parzysty nagłówek
		\large Chapter \thechapter \hspace{16pt} \hrulefill 
		\put(-20, 12.1){%
			\makebox(0,0)[l]{%
				\includegraphics[width=0.05\textwidth]{Figures/wppt-logo}
			}
		}
		\put(-240,5.5){%
			\makebox(0,0)[l]{%
				\small Faculty of Fundamental Problems of Technology
			}
		}
	}
	\fancyfoot[LE,RO]{							% Stopki
		\thepage{}
	}  
	\renewcommand{\headrulewidth}{0pt}			% Grubość linii w nagłówku
	\renewcommand{\footrulewidth}{0.2pt}		% Grubość linii w stopce
}


% Domyślny styl dla bibliografii
\fancypagestyle{bibliographyStyle}{
	\fancyhf{}									% wyczyść stopki i nagłówki
	\fancyhead[RO]{								% Prawy, nieparzysty nagłówek
		\hrulefill \hspace{16pt} \large Bibliography
		\put(-455.4, 13.5){%
			\makebox(0,0)[l]{%
				\includegraphics[width=0.05\textwidth]{Figures/pwr-logo}
			}
		}
		\put(-426,5.5){%
			\makebox(0,0)[l]{%
				\small Wrocław University of Science and Technology
			}
		}
	}
	\fancyhead[LE]{								% Lewy, parzysty nagłówek
		\large Bibliography \hspace{16pt} \hrulefill 
		\put(-20, 12.1){%
			\makebox(0,0)[l]{%
				\includegraphics[width=0.05\textwidth]{Figures/wppt-logo}
			}
		}
		\put(-240,5.5){%
			\makebox(0,0)[l]{%
				\small Faculty of Fundamental Problems of Technology
			}
		}
	}
	\fancyfoot[LE,RO]{							% Stopki
		\thepage
	}
	\renewcommand{\headrulewidth}{0pt}			% Grubość linii w nagłówku
	\renewcommand{\footrulewidth}{0.2pt}		% Grubość linii w stopce
}

% Domyślny styl dla dodatków
\fancypagestyle{appendixStyle}{
	\fancyhf{}									% wyczyść stopki i nagłówki
	\fancyhead[RO]{								% Prawy, nieparzysty nagłówek
		\hrulefill \hspace{16pt} \large Appendix \thechapter
		\put(-455.4, 13.5){%
			\makebox(0,0)[l]{%
				\includegraphics[width=0.05\textwidth]{Figures/pwr-logo}
			}
		}
		\put(-426,5.5){%
			\makebox(0,0)[l]{%
				\small Wrocław University of Science and Technology
			}
		}
	}
	\fancyhead[LE]{								% Lewy, parzysty nagłówek
		\large Appendix \thechapter \hspace{16pt} \hrulefill 
		\put(-20, 12.1){%
			\makebox(0,0)[l]{%
				\includegraphics[width=0.05\textwidth]{Figures/wppt-logo}
			}
		}
		\put(-240,5.5){%
			\makebox(0,0)[l]{%
				\small Faculty of Fundamental Problems of Technology
			}
		}
	}
	\fancyfoot[LE,RO]{							% Stopki
		\thepage{}
	}
	\renewcommand{\headrulewidth}{0pt}			% Grubość linii w nagłówku
	\renewcommand{\footrulewidth}{0.2pt}		% Grubość linii w stopce
}

% Osobny styl dla stron zaczynających rozdział/spis treści itd. (domyślnie formatowane jako "plain")
\fancypagestyle{chapterBeginStyle}{
	\fancyhf{}%
	\fancyfoot[LE,RO]{
		\thepage
	}
	\renewcommand{\headrulewidth}{0pt}
	\renewcommand{\footrulewidth}{0.2pt}
}

% Styl dla pozostałych stron spisu treści
\fancypagestyle{tableOfContentStyle}{
	\fancyhf{}%
	\fancyfoot[LE,RO]{
		\thepage
	}
	\renewcommand{\headrulewidth}{0pt}
	\renewcommand{\footrulewidth}{0.2pt}
}

% Formatowanie tytułów rozdziałów i/lub sekcji
\usepackage{titlesec}

% Formatowanie tytułów rozdziałów
\titleformat{\chapter}[hang]					% kształt
{
	\vspace{-15ex}
	\Huge
	\bfseries
}												% formatowanie tekstu modyfikowanego elementu
{\thechapter}												% etykieta występująca przed tekstem modyfikowanego elementu, niewidoczna w spisie treści
{
	10pt
}												% odstęp formatowanego tytułu od lewego marginesu/etykiety
{
	\Huge
	\bfseries
}												% formatowanie elementów przed modyfikowanym tytułem
[
\vspace{2ex}
% \rule{\textwidth}{0.4pt}
%\vspace{-4ex}
]												% dodatkowe formatowanie stosowane poniżej modyfikowanego tytułu


% Formatowanie tytułów sekcji
\titleformat{\section}[hang]					% kształt
{
	\vspace{2ex}
%	\titlerule\vspace{1ex}
	\Large\bfseries
}												% formatowanie tekstu modyfikowanego elementu
{
	\thesection									% etykieta występująca przed tekstem modyfikowanego elementu, niewidoczna w spisie treści
}
{
	0pt
}												% odstęp formatowanego tytułu od lewego marginesu/etykiety
{
	\Large
	\bfseries
}												% formatowanie elementów przed modyfikowanym tytułem
\title{}
% ***************************************************************************************************
% Linki
% ***************************************************************************************************

% Umożliwia wstawianie hiperłączy do dokumentu
\usepackage{hyperref}							% Aktywuje linki

\hypersetup{
	colorlinks	=	true,					% Koloruje tekst zamiast tworzyć ramki.
	linkcolor		=	blue,					% Kolory: referencji,
    citecolor		=	blue,					% cytowań,
	urlcolor		=	blue					% hiperlinków.
}

% Do stworzenia hiperłączy zostanie użyta ta sama (same) czcionka co dla reszty dokumentu
\urlstyle{same}




% ***************************************************************************************************
% Linki
% ***************************************************************************************************

% Umożliwia zdefiniowanie własnego stylu wyliczeniowego
\usepackage{enumitem}

% Nowa lista numerowana z trzema poziomami
\newlist{myitemize}{itemize}{3}

% Definicja wyglądu znacznika pierwszego poziomu
\setlist[myitemize,1]{
	label		=	\textbullet,
	leftmargin	=	4mm}

% Definicja wyglądu znacznika drugiego poziomu
\setlist[myitemize,2]{
	label		=	$\diamond$,
	leftmargin	=	8mm}

% Definicja wyglądu znacznika trzeciego poziomu
\setlist[myitemize,3]{
	label		=	$\diamond$,
	leftmargin	=	12mm
}


% ****************************** Biblatex ***********************************************************
\usepackage[
    backend=biber,
    style=numeric,
	sorting=none
  ]{biblatex}
\addbibresource{qliom.bib}

% ***************************************************************************************************


% ***************************************************************************************************
% Inne pakiety
% ***************************************************************************************************

% Dołączanie rysunków
\usepackage{graphicx}

% Figury i przypisy
\usepackage{caption}
\usepackage{subcaption}

% Umożliwia tworzenie przypisów wewnątrz środowisk
\usepackage{footnote}

% Umożliwia tworzenie struktur katalogów
\usepackage{dirtree}
% Precyzyjne obliczenia szerokości/wysokości dowolnego fragmentu wygenerowanego przez LaTeX
\usepackage{calc}

% my packages
\usepackage{physics}
\usepackage{microtype}
\usepackage{cancel}
\usepackage{color}
\renewcommand{\CancelColor}{\color{red}}
\usepackage{lipsum}

\usepackage{epigraph}
\setlength\epigraphwidth{11cm}
\setlength\epigraphrule{0pt}

\usepackage{etoolbox}
\makeatletter
\patchcmd{\epigraph}{\@epitext{#1}}{\itshape\@epitext{#1}}{}{}
\makeatother

\usepackage{tikz}
\usetikzlibrary{shapes,arrows,decorations.pathreplacing}

% \usepackage{indentfirst}
\usepackage{bm}
\usepackage{quotchap}
\usepackage{float}

% ***************************************************************************************************
% Matematyczne skróty
% ***************************************************************************************************

% Skrócony symbol liczb rzeczywistych
\newcommand{\RR}{\mathbb{R}}
\newcommand{\CC}{\mathbb{C}}

% Skrócony symbol liczb naturalnych
\newcommand{\NN}{\mathbb{N}}

% Skrócony symbol liczb wymiernych
\newcommand{\QQ}{\mathbb{Q}}

% Skrócony symbol liczb całkowitych
\newcommand{\ZZ}{\mathbb{Z}}

% Skrócony symbol logicznej implikacji
\newcommand{\IMP}{\rightarrow}

% Skrócony symbol  logicznej równoważności
\newcommand{\IFF}{\leftrightarrow}

% Symbole operatorów spinowych
\newcommand{\Id}{\mathbb{1}}
\newcommand{\Sp}{S^{+}}
\newcommand{\Sm}{S^{-}}
\newcommand{\Sz}{S^{z}}
\newcommand{\Sx}{S^{x}}
\newcommand{\Sy}{S^{y}}


\newcommand{\hs}[2]{\left(#1 \middle| #2\right)}
\newcommand{\hsSmall}[2]{(#1 | #2)}
\newcommand{\dimension}{\mathcal{D}}
\newcommand{\ua}{\uparrow}
\newcommand{\da}{\downarrow}

\newcommand{\spaa}{\psi_{\textrm{space}}^{(a)}}
\newcommand{\spas}{\psi_{\textrm{space}}^{(s)}}
\newcommand{\spia}{\psi_{\textrm{spin}}^{(a)}}
\newcommand{\spis}{\psi_{\textrm{spin}}^{(s)}}
\newcommand{\hspi}{\mathcal{H}_{\textrm{spin}}}
\newcommand{\hspac}{\mathcal{H}_{\textrm{space}}}


\DeclarePairedDelimiterX\set[1]\lbrace\rbrace{\def\given{\;\delimsize\vert\;}#1}

\definecolor{wppt}{RGB}{0,156,158}
\definecolor{pinegreen}{RGB}{0,128,128}
\definecolor{pwr}{RGB}{154,52,45}
\renewcommand*{\chapterheadstartvskip}{\vspace*{-1\baselineskip}} 
\renewcommand*{\sectfont}{\bfseries\fontseries{bx}\selectfont}
\makeatletter
\let\size@chapter\huge
\renewcommand*{\chapnumfont}{
    %\usefont{T1}{jkpvos}{b}{n}
    \usefont{T1}{pzc}{b}{n}
    \fontsize{100}{120}\color{wppt}\selectfont}
% default number
% \renewcommand*{\chapnumfont}{%
%   \usefont{T1}{\@defaultcnfont}{b}{n}\fontsize{120}{150}\selectfont% Default: 100/130
%   \color{chaptergrey}}
\makeatother

\newcommand*{\blankpage}{%
\vspace*{\fill}
{\centering This page is intentionally left blank.\par}
\vspace{\fill}}
% \makeatletter
% \renewcommand*{\cleardoublepage}{\clearpage\if@twoside \ifodd\c@page\else
% \blankpage
% \thispagestyle{empty}
% \newpage
% \if@twocolumn\hbox{}\newpage\fi\fi\fi}
% \makeatother

% \makeatletter
% \renewcommand*\cleardoublepage{\clearpage\if@twoside
%   \ifodd\c@page\else\hbox{}\newpage\if@twocolumn\hbox{}%
%   \newpage\fi\fi\fi}
% \makeatother


% ***************************************************************************************************
% Środowiska
% ***************************************************************************************************

% Środowisko do twierdzeń
\newtheorem{theorem}{Theorem}[chapter]

% Środowisko do lematów
\newtheorem{lemma}{Lemma}[chapter]

\newtheorem{proposition}{Proposition}[chapter]

% Środowisko do przykładów
\newtheorem{example}{Example}[chapter]

% Środowisko do wniosków
\newtheorem{corollary}{Corollary}[chapter]

% Środowisko do definicji
\newtheorem{definition}{Definition}[chapter]

% Środowisko do dowodów
\newenvironment{proof}{
	\par\noindent \textbf{Proof.}
}{
\begin{flushright}
	\vspace*{-6mm}\mbox{$\blacksquare $}
\end{flushright}
}

% Środowisko do uwag
\newenvironment{remark}{
	\bigskip \par\noindent \small \textbf{Remark.}
}{
\begin{small}
	\vspace*{4mm}
\end{small}
}

\newenvironment{abstract}[1]{%
	\begin{center}
		\normalfont\textbf{Abstract}
	\end{center}
	\begin{quotation} 
		#1 
	\end{quotation}
}{%
\vspace{1cm}
}


% ***************************************************************************************************
% Dokument
% ***************************************************************************************************

\frontmatter

\begin{document}
\pagenumbering{Roman}
\begin{titlingpage}
	\vspace*{\fill}
	\begin{center}
		\begin{picture}(300,510)
			\put(11,520){\makebox(0,0)[l]{\large \textsc{Faculty of Fundamental Problems of Technology}}}
			\put(11,500){\makebox(0,0)[l]{\large \textsc{Wrocław University of Science and Technology}}}
			% Tytuł pracy
			\put(80,320){\Huge \textsc{Eigenmodes in}}
			\put(80,280){\Huge \textsc{nearly integrable}}
			\put(80,240){\Huge \textsc{quantum chains}}
			% Autor pracy
			\put(90,200){\makebox(0,0)[l]{\large \textsc{Jakub Pawłowski}}}
			\put(90,180){\makebox(0,0)[l]{\large \textsc{Index number: 250193}}}
			
			\put(200,100){\makebox(0,0)[l]{\large Master thesis}}
			\put(200,80){\makebox(0,0)[l]{\large under supervision of}}
			% dane promotora
			\put(200,60){\makebox(0,0)[l]{\large prof.\ dr hab.\ Marcin Mierzejewski}}
			
			\put(115,-70){\includegraphics[width=0.15\textwidth]{Figures/pwr_logo_english.pdf}}
			\put(106,-85){\makebox(0,0)[bl]{\large \textsc{Wrocław 2023}}}
		\end{picture}
	\end{center}
	\vspace*{\fill}
\end{titlingpage}

\pagestyle{tableOfContentStyle}
\blankpage{}
% \cleardoublepage{}
\newpage{}
\pagestyle{tableOfContentStyle}
\vspace*{14cm}
\begin{center}
	\epigraph{\normalsize
		I would like to express my sincere gratitude to prof. dr hab.\ Marcin Mierzejewski for being my
		% guide into the world of quantum many-body physics. Without his immense knowledge and willingness
		% to help, this work would not have been possible.\\
		% \vspace*{0.5cm}
		% I also own my heartfelt appreciation to prof.\ dr hab.\ Katarzyna Sznajd-Weron who
		% sparked my love for theoretical physics and introduced me to academia.\\
		% \vspace*{0.5cm}
		% Last but not least, I would like to thank my parents, for without them I would not be here
		% where I am today.
	}{}
\end{center}
% \cleardoublepage{}
\newpage{}
\blankpage{}
\newpage{}

\pagestyle{tableOfContentStyle}
\vspace*{6cm}
\abstract{\noindent
	TEST
	% unitary quantum dynamics has for a long time been a topic of intensive research.
	% One possible approach to reconciliation between these two fields in generic nonintegrable
	% systems is given by \textit{Eigenstate Thermalization Hypothesis}. On the other hand,
	% in integrable systems thermalization is hindered by the existence of extensive
	% number of (quasi)local integrals of motions, (Q)LIOMs. In this work we investigate a class
	% of \textit{noncommuting} (Q)LIOMs originating from symmetries
	% of spin-\(1/2\) XXZ model present for certain values of anisotropy parameter.
	% However, it is the boundary between chaotic and integrable systems
	% that is the most interesting, as it is still largely unknown. To this end, a weak
	% integrability-breaking perturbation is introduced and the relaxation of aforementioned
	% noncommuting (Q)LIOMs is studied. We report a novel Gaussian-like decay of autocorrelation functions,
	% together with a linear, instead of the expected quadratic, scaling of relaxation times with
	% perturbation strength. This result is a consequence of the origin of these noncommuting
	% quantities being the macroscopical degeneracy of many-body spectra for high-symmetry anisotropy values.
}
\begin{quotation}
	\noindent \textbf{Keywords:} \textit{integrals of motion, ETH, integrability breaking, XXZ model}
\end{quotation}
\newpage{}
\blankpage{}
\newpage{}

\pagestyle{tableOfContentStyle}
\tableofcontents
\newpage{}
\blankpage{}
\newpage{}
% \cleardoublepage{}
% ***************************************************************************************************
% Wstęp
% ***************************************************************************************************

\pagestyle{custom}
\mainmatter{}

% ***************************************************************************************************
% Rodziały
% ***************************************************************************************************

\chapter{Introduction\label{chap:intro}}
\thispagestyle{chapterBeginStyle}

\section{Nearly integrable quantum systems}
One of the most widely disputed problems in modern physics is the reconciliation
of irreversible thermodynamics~\autocite{Huang1987, Feynman2018}
with unitary, time-reversible dynamics predicted by quantum mechanics~\autocite{Landau1976, Sakurai2017}.
In other words, it is whether a generic, isolated quantum system can and will `forget' about its
initial, nonequilibrium state. In recent years, this problem has attracted the attention of many physicists,
especially taking into account experimental evidence of the loss of information, or thermalization, in isolated
systems~\autocite{Trotzky2012, Rigol2012, Rigol2008, Hung2010, Hofferberth2007}. Understanding this phenomenon is a crucial first step to
controlling it, and perhaps to the creation of sought-after, robust systems which do not exhibit such
behavior. After all, `forgetting' about the initial state is equivalent to scrambling of quantum information
and decoherence, which has been known for a long time to be one of the most important hindrances
to a functional and scalable quantum computer~\autocite{Shor1995, Lewis-Swan2019}. Given the multitude
of possible applications of such a device, both purely scientific and commercial, controlling thermalization
in isolated systems could be the next groundbreaking achievement~\autocite{MacQuarrie2020}.

\paragraph{Preventing thermalization} A possible theoretical explanation of the thermalization in isolated quantum systems
was proposed in seminal papers by~\textcite{Deutsch1991,Srednicki1994} in form of an ansatz for matrix elements of
local quantum observables \(A_{mn} = \matrixel{m}{A}{n}\) (in eigenbasis \(\hat{H} \ket{n}=E_n \ket{n}\) of some Hamiltonian),
known as \textbf{Eigenstate Thermalization Hypothesis} (ETH)
\begin{equation}
    A_{mn} = A(\bar{E} ) \delta_{mn} + \mathrm{e}^{-S(\bar{E})/2}f_{A}(\bar{E} ,\omega)R_{mn}
    \label{eq:ETH}
\end{equation}
where \(\bar{E}  = (E_n+E_m)/2\), \(\omega=E_n-E_m\) and \(S(\bar{E} )\) is
the thermodynamic entropy at energy \(\bar{E} \),  \(A(\bar{E} )\) and \(f_{A}(\bar{E} ,\omega)\) are
some smooth functions of their parameters and
\(R_{mn}\) is a random real or complex variable. The origin of ETH is the Random Matrix Theory~\autocite{Wigner1955, Mehta2004},
but it was the key insight of Deutsch and Srednicki that connected it directly to statistical mechanics.
Unraveling the physics behind equation~\eqref{eq:ETH}, we have that, for generic systems,
the long-time (thermalized) expected values of local observables are related to the predictions of Gibbsian ensembles
known from equilibrium statistical physics~\autocite{DAlessio2016}. Diagonal matrix elements \(A_{nn} \simeq A(E_n)\) are smooth functions of energy and
coincide with the microcanonical average at energy \(E_n\), whereas the off-diagonal matrix elements and fluctuations
of the diagonal ones are postulated to decrease exponentially with system size~\autocite{Beugeling2014}.
While ETH was shown to hold analytically in only a few selected systems~\autocite{Magan2016}, numerical clues can be
found in various systems with many-body interactions such as hard-core bosons, interacting spin
chains~\autocite{Santos2010a,Rigol2010,Khatami2013,Rigol2009a} and fermions~\autocite{Neuenhahn2012,Rigol2009}.

Given the lack of a definitive answer to the question of when ETH holds and the fact that our ultimate goal is preventing thermalization, it is desirable to look for systems
that explicitly violate this hypothesis. There are a few known classes of such systems:
quantum integrable models, systems with `scars', which freeze thermalization for
some infrequent, but physically relevant states~\autocite{Turner2018a, Turner2018}, quantum time crystals with
broken time-translation symmetry~\autocite{Wilczek2012, Sacha2017} and systems with many-body interactions together
with some form of disorder~\autocite{Basko2006}. As the thermalization in those systems is slowed down or even
completely stopped, in principle they could preserve information about the initial state for an arbitrarily long time.
This behavior is in stark contrast with what is usually observed in interacting, many-body systems, namely fast
dynamics on the time scales of tens of femtoseconds~\autocite{DalConte2015}.
In this thesis, we shall concern ourselves only with the case of \textbf{quantum integrable models}.

\paragraph{Robustness of integrability}
In quantum mechanics, as opposed to classical mechanics, a unique definition of integrability
is yet to emerge~\autocite{Caux2011, Yuzbashyan2013}. Nevertheless, a common practice is to take as a sufficient condition of
integrability the presence of an extensive number (increasing linearly with the system size) of \textbf{local observables
that are integrals of motion}, i.e., commute with the Hamiltonian, called \textbf{LIOMs} for short. The significance of such
systems is at least twofold.
First, as already mentioned, they can provide a new method of storage of information, encoded in expectation values
of LIOMs, which can pave the way for the creation of the sought-after quantum memory.
Second, they are among a few quantum many-body systems that are susceptible to analytical methods and so provide
a rich playground for both theoretical physicists and mathematicians. Tools such as Algebraic Bethe
Ansatz~\autocite{Faddeev1995,Faddeev1996,Korepin1993} or Generalized Hydrodynamics~\autocite{Agrawal2020,Friedman2020,
    Bertini2021, Bastianello2021, Bulchandani2021} lead to considerable insight regarding the nature of integrable systems.

Unfortunately, it is often the case for integrable systems such as Heisenberg, Hubbard or Lieb-Linger
models whose integrability relies on a certain set of fine-tuned parameters. Any deviation from these parameters, e.g.,
in the form of some perturbation, can very easily destroy the integrability. Exactly this situation takes place in most
experimental setups (albeit some signatures of integrability were observed~\autocite{Khemani2019}, suggesting sufficient
proximity to an integrable system),  even though the capabilities and precision of control have
reached remarkable levels~\autocite{Browaeys2020}. Therefore, more realistic systems
are expected to be described by \textbf{nearly integrable systems}, containing some non-negligible perturbations that
impact integrability~\autocite{Zotos2004, Brandino2015}.
This renders all the formalism for investigations of purely
integrable systems inapplicable and forces one to rely mostly on numerical methods (for a recent review
see~\textcite{Bertini2021}).
A general expectation for such nearly integrable systems is such, that in the thermodynamic limit, an
arbitrarily small perturbation should restore generic chaotic dynamics, which leads to thermalization
~\autocite{LeBlond2021}. However, there
are some reports about surviving traces of integrability, for example in the form of residual quasiconserved
quantities~\autocite{Brandino2015}. In most cases numerical methods can only access finite
systems and indirect transitions to infinite system sizes, such as finite-size scaling, pose great difficulties
when executed properly, e.g., first thermodynamic limit and then infinite time limit~\autocite{Sirker2014, CamposVenuti2010}.
Therefore, it is important to first gain a deep understanding of \textbf{weakly-perturbed finite systems}. Such
endeavor immediately raises a few questions: How strong should be a perturbation in a finite system
to destroy integrability, and how to efficiently describe resulting slow dynamics?
Recently, those and other similar questions have been asked concerning the phenomena of ergodicity breaking
phase transitions~\autocite{Suntajs2020,Suntajs2022}.
Insight from experiments with ultracold atomic gases suggest also another way for nearly integrable systems to emerge,
namely generic, chaotic one- and quasi-one-dimensional systems that exhibit approximately integrable dynamics on experimentally relevant timescales,
because of a lack of rapid local thermalization~\autocite{Polkovnikov2011, Gopalakrishnan2023}.

In classical mechanics, answers to those types of questions are given by the Hamiltonian perturbation
theory~\autocite{Arnold1990}. There, the distinction between integrable and non-integrable systems is
well understood. It is known, that for classical systems with \(n\) degrees of freedom
to be integrable, it is sufficient to have \(n\) integrals of motion \(\{H,F_i\} = 0,\; i \in \{1,\ldots,n\} \) that
are in involution, i.e., \(\{F_i,F_j\} = 0\) for any \(i,j\in\{1,\ldots,n\}\). Then, dynamics are restricted to an \(n\)-dimensional
torus in the phase space. The fate of such invariant tori under integrability-breaking perturbations is given by the
famous Kolmogorov-Arnold-Moser theorem (KAM)~\autocite{Kolmogorov1954, Arnold1963, Moser1962}, stating that for systems
with finite degrees of freedom, the majority of the invariant tori occupying the phase space survive the influence of small
perturbations. However, as of now, there is no equivalent of this result in quantum mechanics.
Nevertheless, such quantum systems are very intriguing as they can facilitate robust prethermalization plateaux,
i.e., dynamics that at intermediate times resemble that of integrable models
(at least for suitably weak perturbations), even though the system eventually thermalizes at longer times~\autocite{Mallayya2019, Bertini2015,Berges2004,Langen2016}.

\section{Motivation and aim of this dissertation}
As argued in the previous section, nearly integrable quantum systems constitute an important relaxation
of constraints imposed on ordinary integrable systems, facilitating experimental realizations. They form the broader context this
thesis is set in and the main motivation. However, to keep it finite in size, we narrowed
our scope to only one concrete physical problem — the dynamics of spin transport in the long-range
anisotropic Heisenberg spin chain.

For completeness, the initial part of this thesis is a pedagogical review of the so-called
\textbf{Krylov subspace methods} which allow us to leverage the sparsity of local observables
in order to avoid the limitations originating from the exponential growth of many-body Hilbert space.
Following the exposition by~\textcite{Trefethen1997}, in Chapter~\ref{chap:krylov} we will start our journey seemingly far away from
physics, investigating a general algorithm called the Arnoldi iteration, which original aim was to reduce
a matrix to an `almost triangular' or Hessenberg form. Along the way, we will discover that, as a byproduct,
Arnoldi iteration produces an excellent approximation of extremal eigenvalues and eigenvectors. Contrary
to most physics texts and in the spirit of the pedagogical nature of this chapter, we will try to derive all results
when possible and motivate them thoroughly when not. Setting course back to physics, next, we will assume that our matrices are
Hermitian and observe how the Arnoldi iteration simplifies tremendously, producing the well known
Lanczos iteration~\autocite{Sandvik2010}. Then, going beyond just groundstate calculations, we will show how to slightly
modify Lanczos iteration to be able to compute an approximation of action of any analytic function of a Hermitian
matrix on a vector. Applying this scheme to the function \(f(x) = \exp\left(-i x t\right)\), we will obtain an efficient
way of calculating pure state time evolution, called the Krylov propagator~\autocite{Park1986}.
At the end of Chapter~\ref{chap:krylov}, we will outline the recent concept of Quantum Typicality~\autocite{Bartsch2009}
and derive a numerical procedure for efficient calculation of correlation functions, without the need for Exact
Diagonalization. All algorithms described in chapter~\ref{chap:krylov} were, for the purpose of calculations
in this thesis, implemented using
the Armadillo linear algebra library~\autocite{Sanderson2016} and Intel MKL.

The main part of this manuscript is devoted to demonstrating, that for specific parameters, the
 \textbf{long-range anisotropic Heisenberg model} is a nearly integrable system.
Existence of many interesting features of quantum many-body systems, such as ballistic dynamics in Heisenberg spin
chains~\autocite{Zotos1997,Bertini2021},
exotic frustrated magnetism~\autocite{Nisoli2017}, peculiar phase transitions~\autocite{Sandvik2010a,Yang2021}
and entangled spin liquids~\autocite{Balents2010} depends strongly
on the type and range of interaction present in them, so it is safe to say that it plays a crucial role.
In the last few years, considerable development of experimental techniques has
occurred, allowing for unprecedented manipulation of the interactions.
Platforms such as individually controlled Rydberg atoms, or optical lattices provide insight into the
properties of quantum many-body systems. Whereas optical lattice facilitates mostly fermionic
systems~\autocite{Bakr2009, Greif2016, Parsons2015, Boll2016}, Rydberg atoms can be used to simulate pure spin systems.
Models such as Ising or XY emerge naturally from their properties~\autocite{Browaeys2020}, and the ability to control
the range of interactions make them suitable for the simulation of long-range models~\autocite{Borish2020}. Using time-periodic
driving one can turn naturally existing Hamiltonian into some other, effective one - the so-called Floquet Hamiltonian.
So far, this method has been applied with great success to create tunable XXZ model~\autocite{Scholl2022}, strongly distance
selective interactions~\autocite{Hollerith2022}, and tunable XYZ models~\autocite{Steinert2022, Geier2021}.
Progress in experimental methods sparked theoretical interest in long-range
models~\autocite{Richerme2014,Jurcevic2014,Hauke2013,Foss-Feig2015,Maghrebi2016,Lepori2017,Frerot2017,Vanderstraeten2018,Cevolani2018,Kloss2019,Ren2020,Bulchandani2022a},
yet their dynamical properties are still largely unknown, which explains our interest in long-range anisotropic
Heisenberg model in this thesis.
In Chapter~\ref{chap:spin_transport}, motivated by experiments with density expansion in cold atoms~\autocite{Ronzheimer2013,Vidmar2013,Neyenhuis2017},
we will study the dynamics of spin domains using the Krylov propagator, followed by Exact Diagonalization studies of the optical conductivity.
Chapter~\ref{chap:currents} is concerned with investigating a class of local observables exhibiting similar properties
to the spin current, using a numerical procedure searching for most conserved operators~\autocite{Mierzejewski2015a}.
It is important to stress, that the results presented in Chapter~\ref{chap:spin_transport} were not obtained
by the author of this thesis, but by other co-authors. Nevertheless, they are included here in order to ensure
the coherence of the narrative. On the other hand, the results presented in Chapter~\ref{chap:currents} were
obtained solely by the author of this thesis. Both chapters~\ref{chap:spin_transport} and~\ref{chap:currents}
contain \textbf{original results, first published in~\textcite{Mierzejewski2023}.} 

In the remainder of this chapter, we will provide a short introduction to the long-range anisotropic Heisenberg model
and other quantities of interest in this thesis.

\section{Long range anisotropic Heisenberg model and spin current}

We study the paradigmatic quantum model of magnetism, the anisotropic Heisenberg
model, but enriched with a long-range exchange \(J(r) = J/r^{\alpha}\). The full Hamiltonian,
defined on a one-dimensional lattice with periodic boundary conditions, reads
\begin{equation}
    \hat{H}  = \sum_{\ell=1}^{L}\sum_{r=1}^{r_{\mathrm{max}}} J(r) \left[\frac{1}{2} \left(
        \Sp_{\ell}\Sm_{\ell+r} + \Sm_{\ell}\Sp_{\ell+r}\right) + \Delta \Sz_{\ell}\Sz_{\ell+r}
        \right]
    \label{eq:long_range}
\end{equation}
where \(r_{\mathrm{max}}\) is taken to be \(\left\lceil L/2\right\rceil - 1 \), to avoid double counting of hoppings.
Unless stated otherwise, we will work in units where \(J = 1\).
The spin\(-\frac{1}{2}\) operators are defined in the usual way
\begin{equation}
    S_{\ell}^{a} \equiv  \underbrace{\Id \otimes \Id \otimes \ldots \Id}_{\ell-1} \otimes S^{a} \otimes \underbrace{\Id \otimes \ldots \otimes \Id}_{L-\ell}
    \label{eq:spin_op}
\end{equation}
where \(a\in\{+,-,z\}\), \(\Sp = S^x + i S^y\), \(\Sm = \left(\Sp\right)^{\dagger}\),
\(\Id\) is a \(2\cross 2\) identity matrix and operators \(S^x,S^y,S^z\) are defined in terms
of corresponding Pauli matrices and obey the \(\mathfrak{su}(2)\) algebra commutation relations
\begin{equation}
    \comm{S^a}{S^b} = i \varepsilon_{abc} S^c
\end{equation}
From equation~\eqref{eq:spin_op} it is evident, that the Hilbert space \(\mathcal{H}\) our Hamiltonian acts on
is the tensor product of \(L\) copies of single spin spaces \(\mathfrak{h}_{\ell} \cong \CC^2\),
\begin{equation}
    \mathcal{H} = \bigotimes_{\ell=1}^{L} \mathfrak{h}_{\ell}
\end{equation}
As \(\mathrm{dim_{\CC}}(\mathfrak{h}) = 2\), the total dimension is \(2^L\).
For numerical calculations, we shall use the so-called \textbf{Ising basis}, consisting of eigenstates
of the \(S^z_{\mathrm{tot}} = \sum_{\ell=1}^{L} \Sz_{\ell}\) operator. As on each site we can only
have spin up or down, an efficient representation of such states is achieved using binary numbers.

Our main interest is in the spin transport in this system. To this end, we first notice that
the total magnetization, or \(z\) component of the total spin \(\Sz_{\mathrm{tot}}\)
is a conserved quantity.
It can be shown by directly evaluating the commutator \(\comm{\hat{H} }{\Sz_{\mathrm{tot}}}=0\), or by just observing
that both spin-flip and interaction terms do not change the number of up/down spins when acting on a state
from the Ising basis. Thus, we can write down a well-defined continuity equation for the evolution
of local spin density
\begin{equation}
    \dv{t} \Sz_{\ell}(t) + \div{j^{\sigma}_{\ell}(t)} = 0
    \label{eq:spin_cont_eq}
\end{equation}
where \(\div{j^{\sigma}_{\ell}(t)} = j^{\sigma}_{\ell+1}(t) - j^{\sigma}_{\ell}\) is the
discrete divergence of the spin current and \(\Sz_{\ell}(t) = e^{i \hat{H}  t} \Sz_{\ell} e^{-i \hat{H}  t}\).
The time derivative in the above equation is given by the Heisenberg equation
\begin{equation}
    \dv{t} \Sz_{\ell}(t) = i \comm{\hat{H} }{\Sz_{\ell}(t)}
    \label{eq:Heisenberg}
\end{equation}
Combining equations~\eqref{eq:spin_cont_eq} and~\eqref{eq:Heisenberg}, we get
\begin{equation}
    j^{\sigma}_{\ell+1} - j^{\sigma}_{\ell} = i\comm{\Sz_{\ell}}{\hat{H} }
\end{equation}
This equation is simple in principle, however, its solution can get tedious in case of more complicated
operator densities.

Fortunately, there is a trick that yields a simple expression for the current
associated with an arbitrary conserved, extensive operator. Let \(X = \sum^{L}_{\ell=1}x_l\) be such quantity,
with \(x_{\ell}\) being its local density and \(j^x = \sum^{L}_{\ell=1} j^x_{\ell}\) being the current.
Then define the polarization operator
\begin{equation}
    P \equiv  \sum^{L}_{\ell=1}\ell x_{\ell} 
\end{equation}
We will now calculate its derivative in two ways, first using the continuity equation and second using
the Heisenberg picture. For brevity, we suppress the explicit time dependence of operators.
\begin{align}
    \dv{P}{t} & = \sum^{L}_{\ell=1} \ell \dv{x_{\ell}}{t} =
    \sum^{L}_{\ell=1} \ell j^x_{\ell} - \underbrace{\sum_{\ell=1}^{L} \ell j^x_{\ell+1}}_{\ell \to \ell - 1}
    = \sum_{\ell=1}^{L} \left[\ell j^x_{\ell} - \left(\ell-1\right)j^{x}_{\ell}\right]
    = \sum_{\ell=1}^L j^x_{\ell} = j^x                                                    \\
    \dv{P}{t} & = \dv{t} \sum_{\ell=1}^L \ell \left(e^{i \hat{H}  t}x_{\ell} e^{-i \hat{H}  t}\right) =
    \sum_{\ell}^{L} \left(\ell e^{i \hat{H}  t} i \comm{\hat{H} }{x_{\ell}}e^{-i \hat{H}  t}\right) =
    e^{i \hat{H}  t} \left(i \sum_{\ell} \ell \comm{H}{x_{\ell}}\right)e^{-i \hat{H}  t}
\end{align}
From the two equations above we can read off a compact expression for the current operator in the Schr{\"o}dinger
picture
\begin{equation}
    j^x = i \sum_{\ell=1}^L \ell\comm{\hat{H} }{x_{\ell}}
    \label{eq:polarization}
\end{equation}
Substituting \(x_{\ell} = \Sz_{\ell}\) we now quickly derive
\begin{align}
    j^{\sigma} & = i \sum_{\ell=1}^L \comm{\hat{H} }{\Sz_{\ell}} = i \sum_{\ell,\ell^{\prime}=1}^L \sum_{r=1}^{r_{\mathrm{max}}}
    \ell J(r) \left(\frac{1}{2}\comm{\Sp_{\ell^{\prime}}\Sm_{\ell^{\prime}+r}}{\Sz_{\ell}}
    + \frac{1}{2}\comm{\Sm_{\ell^{\prime}}\Sp_{\ell^{\prime}+r}}{\Sz_{\ell}}
    + \Delta \cancel{\comm{\Sz_{\ell^{\prime}}\Sz_{\ell^{\prime}+1}}{\Sz_{\ell}}} \right) \nonumber                       \\
               & = \frac{i}{2}  \sum_{\ell,\ell^{\prime}=1}^L \sum_{r=1}^{r_{\mathrm{max}}}\ell J(r)  \left(
    \delta_{\ell,\ell^{\prime}+r}\Sp_{\ell^{\prime}} \Sm_{\ell^{\prime}+r}  - \delta_{\ell,\ell^{\prime}}
    \Sp_{\ell}\Sm_{\ell^{\prime}+r} - \delta_{\ell^{\prime} + r, \ell} \Sm_{\ell^{\prime}} \Sp_{\ell} +
    \delta_{\ell^{\prime},\ell}\Sm_{\ell}\Sp_{\ell^{\prime}+r}
    \right) \nonumber                                                                                                     \\
               & = \frac{i}{2}  \sum_{\ell^{\prime}=1}^L \sum_{r=1}^{r_{\mathrm{max}}} J(r)  \left(
    (\ell^{\prime}+r)\Sp_{\ell^{\prime}} \Sm_{\ell^{\prime}+r}  - \ell^{\prime}
    \Sp_{\ell^{\prime}}\Sm_{\ell^{\prime}+r} - (\ell^{\prime}+r) \Sm_{\ell^{\prime}} \Sp_{\ell^{\prime}} +
    \ell^{\prime}\Sm_{\ell^{\prime}}\Sp_{\ell^{\prime}+r}
    \right) \nonumber                                                                                                     \\
               & = \frac{i}{2} \sum_{\ell=1}^L \sum_{r=1}^{r_{\mathrm{max}}} \frac{J}{r^{\alpha-1}}
    \left(\Sp_{\ell}\Sm_{\ell+r}-\Sm_{\ell}\Sp_{\ell+r}\right)
    \label{eq:spin_current}
\end{align}
which is our desired spin current. It will be the quantity of central interest in Chapter~\ref{chap:spin_transport}.

Before we finish the introductory chapter, let us notice two symmetries shared by the
Hamiltonian~\eqref{eq:long_range} and spin current~\eqref{eq:spin_current}, which are particularly useful for
numerical calculations. First of them, we have just met -- it manifests itself as the conservation of magnetization,
i.e. \(\comm{\hat{H} }{\Sz_{\mathrm{tot}}}=0\).
This \(U(1)\) symmetry allows us to decompose the full Hilbert space into parts consisting of states with the same total \(z\)-component
of spin. In more mathematical terms, we have the following
\begin{equation*}
    \mathcal{H} = \bigoplus_{j = 0}^{L} \mathcal{H}_j \text{, where } \big(\forall \ket{\psi} \in \mathcal{H}_j \big) \big(\Sz_{tot} \ket{\psi} = \frac{1}{2}(2j-L) \ket{\psi} \big)
\end{equation*}
i.e., the full Hilbert space with \(\dim{\mathcal{H}} = 2^L\) can be decomposed into the direct sum of its subspaces
\(\mathcal{H}_j\) such that \(\dim{\mathcal{H}_j} = \binom{L}{j}\) and all states in a given subspace correspond to the same
eigenvalue of \(\Sz_{\mathrm{tot}}\) operator. The index \(j\) denotes the number of sites with spin up.
It turns out that the Ising basis we use as a default for numerical calculations in spin systems
is already the `correct' basis, as it is enough to just sort the states according to the number of
spins pointing up. This yields a manifestly block diagonal structure of the Hamiltonian matrix
(cf. Figure~\ref{fig:sparse_structure} with a matrix of nearest neighbors Heisenberg model).

Because we have assumed periodic boundary conditions, both the Hamiltonian and the spin current
are invariant under translations by any number of sites - that is under the action of \(\ZZ_L\) cyclic group.
This symmetry allows us to further decompose the Hilbert space
using the eigenbasis of \(T\), the one-site translation operator defined on the Ising basis as
\begin{equation}
    T \ket{\sigma_1,\sigma_2,\ldots,\sigma_L} = \ket{\sigma_L,\sigma_1,\ldots,\sigma_{L-1}}
\end{equation}
where \(\sigma_{\ell} \in \{\downarrow, \uparrow\}\). It is easy to see that \(T\) is unitary,
thus its eigenvalues must lie on the unit circle in the complex plane~\autocite{Horn2012}.
In a finite chain of length \(L\) we have \(T^L = \Id_{L\cross L}\) and the eigenvalues of \(T\) are
quantized, i.e. of the form \(\mathrm{e}^{i k}\), where \(k \in \left\{0,\frac{2\pi}{L},\ldots,\frac{2\pi}{L}(L-1)\right\}\).
As a consequence, the Hilbert space can be decomposed into the direct sum of eigenspaces of \(T\)
\begin{equation}
    \mathcal{H} = \bigoplus_{m=0}^{L-1} \mathcal{H}^{\prime}_m \text{, where } \big(\forall \ket{\psi} \in \mathcal{H}^{\prime}_m \big) \big(T \ket{\psi} = \exp\left({i \frac{2\pi}{L}m}\right) \ket{\psi}\big)
\end{equation}
and once again the Hamiltonian and the spin current are block diagonal in this basis, which each block being
about \(L\) times smaller than the full Hilbert space.
The practical construction of the momentum states is more complicated than that of the magnetization states,
although is not that difficult. We refer the reader to~\textcite{Sandvik2010} for a beautiful and detailed explanation.
Importantly, we also have \(\comm{T}{\Sz_{\mathrm{tot}}}=0\) and thus the two symmetries are compatible.
Using both of them at the same time, the size of the largest block is reduced to \(\binom{L}{L/2}\frac{1}{L}\), which is
still exponentially large in \(L\), but much smaller than the full Hilbert space of dimension \(2^L\).

To end this chapter we note that these two symmetries are not the only ones present in this model.
For a comprehensive treatment of this topic, together with details of numerical implementation, we 
once again refer the reader to the standard text on this topic~\autocite{Sandvik2010}.

\cleardoublepage{}

\chapter{Summary}
\thispagestyle{chapterBeginStyle}


The overarching goal of the work presented in this thesis, serving as the background, is to understand the properties of
quantum systems that are close to integrability, as they facilitate real-world implementation, while
still maintaining some desired properties of strictly integrable systems.
To aid this development, there is a strong need for a robust set of theoretical tools that one
could use for the investigation of such systems. As most of the sophisticated machinery developed
for integrable systems cease to work for the case of nearly integrable systems, one has
to resort to numerical methods. It is impossible to touch upon all the aspects of this
topic in a short thesis, thus we had two concrete goals in mind.

We started with a comprehensive introduction to the numerical methods, beyond the simplest Exact
Diagonalization, that are used in the study of quantum many-body systems. To this end, we presented
the Krylov subspace methods, which are designed to leverage the sparsity of matrix representations of
physical observables, in order to speed up the computations. Starting from the very beginning, we
introduced in detail the Arnoldi iteration, which is most often used to find the extremal eigenvalues
of general, non-Hermitian matrices, even though it does so rather accidentally. We then moved on to
the case of Hermitian matrices, where this procedure reduces to the well-known and versatile
Lanczos iteration. As an immediate application of the Lanczos iteration, beyond just the ground state
computations, the so-called Krylov propagator was presented, which allows one to compute the time evolution of
a state vector, without the need for diagonalization of the Hamiltonian. Finally, we described,
and partially derived, an approach to the correlation functions, based on the Krylov
propagator and the concept of Quantum Typicality, which allows one to compute them
without explicitly carrying out the trace over all many-body states.

The main goal was to apply those methods to the study of the dynamics of the spin transport in the anisotropic
Heisenberg model with power-law interaction \(J(\alpha) = J/r^{\alpha}\), which is an important example of a
long-range interacting system. Motivated by experimental setups, we started by considering the
expansion of a domain-wall initial state and found a non-monotonic dependence of the
center of mass velocity on the anisotropy parameter. By considering maximal velocities,
we found a line in the parameter space, \(\Delta_O(\alpha) = \exp(-\alpha + 2)\),
for which the velocity of spin domain wall expansion is similar to the case of free
nearest-neighbors particles, indicating short-time (high-frequency) ballistic transport.
Next, we studied spin transport from the perspective of the linear response theory,
by considering the optical conductivity and its integral over a frequency window.
It revealed a regime of quasi-ballistic transport, where the spectral sum
of the optical conductivity accumulated at frequencies below \(\Omega^{\ast}/J \sim
10^{-3}-10^{-2}\). As a result, along the line of optimal anisotropy, the spin
transport is ballistic for surprisingly long times, up to \(t \sim 1/\Omega^{\ast}\),
which is in agreement with the results of the domain-wall expansion. The entirety of
the results presented in Chapter~\ref{chap:spin_transport} were obtained by the coauthors of~\textcite{Mierzejewski2023}.
We also investigated the properties of local, slowly relaxing observables
(LSROs), obtained using a numerical algorithm designed to find local integrals of
motion in integrable systems. For completeness, we provided a rather detailed
description of the algorithm itself, together with a possible upgrade, replacing
ED-based time averaging with Lanczos-based long-time correlation functions.
Unfortunately, the upgraded approach turned out to be unsuitable for the purpose of
this thesis, as matrices of long-range operators are no longer sparse and thus do not
benefit from the Lanczos iteration. Nevertheless, using the original algorithm, we
found that the best LSROs decay the slowest precisely for the optimal anisotropy, obtained from
previous methods. Additionally, we showed that they can be understood in terms of projections onto
current-like operators, corresponding via the Jordan-Wigner transformation to long-range fermionic hoppings,
even though the Hamiltonian does not admit a representation in terms of two-body fermionic operators.
Finally, using the finite-time averaging version of the algorithm, we investigated the frequency dependence
of the LSROs, akin to the optical conductivity, and found that the concentration of the spectral weight
resembles the one of the optical conductivity. Chapter~\ref{chap:currents} contains the original work
of the author of this manuscript.

The main result of this thesis, and the broader research published in \textcite{Mierzejewski2023}, is the
discovery of the line of optimal anisotropy, \(\Delta_{\mathrm{O}}(\alpha) = \exp(-\alpha +2)\),
along which the spin transport is ballistic for surprisingly long times \(tJ \sim 10^{-2} -10^{-3}\).
This line smoothly interpolates between integrable two models with purely ballistic
transport, namely free particles for \(\alpha =\infty,\; \Delta =0\) and Haldane-Shastry
like model for \(\alpha  = 2,\; \Delta  = 1\). It suggests that the long-range Heisenberg
model can be thought of as nearly-integrable far beyond just the vicinity of the
integrable limits, but rather for a whole line of parameters. Unfortunately, the precise
microscopic origin of the optimal anisotropy is still unknown. A possible explanation
would be that for \(\delta = \exp(-\alpha  + 2)\) this model is very close to some,
more complicated, integrable model, in which the spin current is an integral of motion.
Investigation of systems similar to the one studied in this thesis, from the perspective
of the search for remnants of integrability, is a promising direction for future research.

\cleardoublepage{}

%%%%%%%%%%%%%%%%%%%%%%%%%%%%%%%%%%%%%%%%%%%%%%%%%%%%%%%%%%%%%%%%%%%%%%%%%%%%%%
%%%%%%%%%%%%%%%%%%%%%%%%%%%%%%% BIBLIOGRAFIA %%%%%%%%%%%%%%%%%%%%%%%%%%%%%%%%%
%%%%%%%%%%%%%%%%%%%%%%%%%%%%%%%%%%%%%%%%%%%%%%%%%%%%%%%%%%%%%%%%%%%%%%%%%%%%%%

\pagestyle{bibliographyStyle}
\printbibliography{}
\thispagestyle{chapterBeginStyle}
\addcontentsline{toc}{chapter}{Bibliography}
\cleardoublepage{}

%%%%%%%%%%%%%%%%%%%%%%%%%%%%%%%%%%%%%%%%%%%%%%%%%%%%%%%%%%%%%%%%%%%%%%%%%%%%%%
%%%%%%%%%%%%%%%%%%%%%%%%%%%%%%%%% DODATKI %%%%%%%%%%%%%%%%%%%%%%%%%%%%%%%%%%%%
%%%%%%%%%%%%%%%%%%%%%%%%%%%%%%%%%%%%%%%%%%%%%%%%%%%%%%%%%%%%%%%%%%%%%%%%%%%%%%

\appendix
\pagestyle{appendixStyle}
\chapter{Appendix\label{app:1}}
\thispagestyle{chapterBeginStyle}

\cleardoublepage{}

\end{document}

